During 700\,$\mu$s, the radio wave travels about
$r = ct = 3\times10^{10}\times700\times10^{-6} = 2.1\times10^{7}$\,cm. The
region from where the burst originates must be no larger than the distance
that light can travel during the duration of the burst. So we estimate that
$r$ is the size of the source. ($r = 2.1\times10^{7}$\,cm,
$R = 2.3$\,kpc) \hfill(10 points)

Flux density observed at 1660\,MHz is
\begin{align*}
S_{\mathrm{1660\,MHz}} &= 0.35\ \mathrm{kJy}\\
  &= 0.35\times10^{3}\times10^{-23}\
     \mathrm{ergs\ s^{-1}\ cm^{-2}\ Hz^{-1}}\\
  &= 3.5\times10^{-21}\ \mathrm{ergs\ s^{-1}\ cm^{-2}\ Hz^{-1}}
\end{align*}
\hfill(10 points)

The solid angle subtended by this source of radiation is
\begin{align*}
\Omega = \pi\,(r/R)^2
  &= 3.14\,\bigl(2.1\times10^{7}/2.3\times10^{3}\times
     3.086\times10^{18}\bigr)^2\\
  &= 2.75\times10^{-29}\ \mathrm{sr}
\end{align*}
\hfill(30 points)

The flux density is related to the total brightness by the relation:
\[ S_{\mathrm{1660\,MHz}} = B_{\mathrm{1660\,MHz}}\,\Omega \]
\hfill(20 points)

while at this frequency, the total brightness can be approximated from the
Rayleigh--Jeans formula:
\[ B_{\mathrm{1660\,MHz}} = 2kT_b\nu^2/c^2 \]
\hfill(20 points)

where $T_b$ is the brightness temperature. Then we have:
\begin{align*}
T_b &= S_{\mathrm{1660\,MHz}}c^2/2k\nu^2\Omega\\
T_b &= [3.5\times10^{-21}\times(3\times10^{10})^2]/
       [2\times(1.38\times10^{-16})\times(1660\times10^{6})^2
        \times(2.75\times10^{-29})]\\
    &= 1.5\times10^{26}\ \mathrm{K}
\end{align*}
\hfill(10 points)
